\subsection{一般情形}

在本节中，G 表示一个局部紧交换群。

命题 3.　a) 存在一个整数 \(n \geqslant 0\) 及一个子群 L，使得 G 是 L 与一个同构于 \(\mathbf{R}^n\) 的子群的直积，而且 L 具有一个紧开子群 K，使得 L/K 是离散的；

\begin{enumerate}
\def\labelenumi{\alph{enumi})}
\setcounter{enumi}{1}
\tightlist
\item
  群 G 是开子群的一个递增滤过族的并，其中每个开子群都是同构于形如 \(R^{p} \times T^{q} \times Z^{r} \times A\) 的群的群的射影极限，这里 A 是有限群，p、q、r 是正整数。
\end{enumerate}

我们证明 \(b\) )。对 \(e\) 的每个紧邻域 V，用 \(\mathrm{G_V}\) 表示由 V 生成的 G 的子群。它是开的，且由命题 2 的推论，群 \(\mathrm{G_V}\) 是形如 \(\mathbf{R}^p\times \mathbf{T}^q\times \mathbf{Z}^r\times \mathrm{A}\) 的群的射影极限，其中 A 是紧群，\(p,q\) 与 \(r\) 属于 \(\mathbf{N}\)。当 V 取遍 \(e\) 的紧邻域时，这些子群 \(\mathrm{G_V}\) 构成一个滤过族（因为对 e 的任意紧邻域 V 与 W，\(G_{V}\) 与 \(G_{W}\) 都包含于 \(G_{V\cup W}\)）。最后，群 G 是诸子群 \(G_{V}\) 的并。

设 H 是 G 的由 e 的一个紧邻域生成的开子群。存在一个紧群 K 及正整数 p 和 q，使得 H 同构于 \(R^{p} \times Z^{q} \times K\)（II，第 246 页，命题 1）；把 H 与这个乘积等同。H 到可除群 \(R^{p}\) 上的典范满射可以延拓为一个态射 \(\pi\)（G 到 \(R^{p}\) 上的态射）（II，第 245 页，引理 3）。它是一个投影子 \(\pi\)（G 到 \(R^{p}\) 上的），由于它在开子群 H 上的限制是连续的，它是连续的。于是 G 是 \(R^{p}\) 与 \(\pi\) 的核 L 的直积（TG，III，第 47 页，推论）。我们有 \(Z^{q} \times K = H \cap L\)，因而 \(Z^{q} \times K\) 是 L 的一个开子群。这样，K 是 L 的一个紧开子群，从而 L/K 是离散的。

命题 4.　设 B \(_{G}\) 是 G 中生成 G 的一个相对紧子群的那些元素组成的集合。那么 B \(_{G}\) 是 G 的一个闭子群，且 B \(_{G}^{\perp}\) 是 \(\widehat{G}\) 的单位元连通分支。

集合 \(B_{G}\) 是 G 的一个子群，因为 G 的两个紧子集的乘积是 G 的一个紧子集。

设 H 是 G 的由 e 的一个紧邻域生成的开子群。存在正整数 p 和 q 及一个紧群 K，使得子群 H 同构于 \(R^{p} \times Z^{q} \times K\)（II，第 246 页，命题 1）。若把这些群等同起来，就看到 \(B_{G} \cap H = K\) 在 H 中是闭的。由于由紧邻域生成的开子群 H 组成的族是 G 的一个开覆盖（例如，对 e 的每个固定的紧邻域 U，x 都属于由 \(U \cap \{x\}\) 生成的子群），由此推出 \(B_{G}\) 是闭的。

现在来计算 \(B_{G}^{\perp}\)。由命题 3 的 a)，存在一个正整数 \(n \geqslant 0\) 及一个具有紧开子群的群 L，使得 G 可以与 \(R^{n} \times L\) 等同。于是 \(B_{G}\) 与 \(\{0\} \times B_{L}\) 等同，而 \(B_{G}^{\perp}\) 与 \(R^{n} \times B_{L}^{\perp}\) 等同。这样我们就归结到 G = L 具有一个紧开子群 K 的情形。

这时有 \(K \subset B_{G}\)。若把 \(\widehat{G/K}\) 与 \(K^{\perp}\) 等同（II，第 226 页，定理 4），则正交 \((\mathrm{B}_{\mathrm{G}}/\mathrm{K})^{\perp}\)（即 \(B_{G}/K\) 在 \(\widehat{G/K}\) 中的正交）与 \(B_{G}^{\perp}\) 等同。但另一方面 \(B_{G}/K = B_{G/K}\)，且由于 \(K^{\perp}\) 是 \(\widehat{G}\) 的一个开子群（II，第 233 页，推论 2），\(\widehat{G}_{0}\)（\(\widehat{G}\) 的中性分支）也是 \(K^{\perp}\) 的中性分支。于是关于离散群 G/K 的断言与关于 G 的断言等价。

最后，假设 G 是离散的。这时群 \(\widehat{G}\) 是紧的（II，第 233 页，命题 18）。中性分支 \((\widehat{\mathrm{G}})_{0}\) 是 \(\widehat{G}\) 的诸开子群的交（TG，III，第 35 页，命题 14），而 \(\widehat{G}\) 的一个子群是开的，当且仅当它是闭的且具有有限指标，或者等价地，它的正交是有限的（II，第 233 页，推论 2）；II，第 228 页的推论 4 表明 \((\widehat{\mathrm{G}})_{0}^{\perp}\) 是 G 的有限子群的并，由于 G 是离散的，它不是别的，正是 \(B_{G}\)。我们由对偶性得出结论。

推论 1.　假设 G 是紧的。那么下列条件等价：

\begin{enumerate}
\def\labelenumi{(\roman{enumi})}
\item
  群 G 是连通的；
\item
  群 \(\widehat{\mathbf{G}}\) 是无挠的；
\item
  群 G 是可除的。
\end{enumerate}

由于 \(B_{G} = G\)，由命题 4 及 II，第 229 页的推论 7 即得。

推论 2.　假设 G 是紧的。那么 G 是全不连通的，当且仅当 \(\widehat{\mathbf{G}}\) 是一个挠群。

群 G 是全不连通的，当且仅当它的单位元连通分支退化为 \(\{e\}\)；命题 4 表明这一条件等价于 \(B_{\widehat{G}} = \widehat{G}\)。由于群 \(\widehat{G}\) 是离散的（II，第 233 页，命题 18），这就相当于说 \(\widehat{G}\) 的每个元素都生成一个有限群，即 \(\widehat{G}\) 是挠群。

推论 3.　若 G 是连通的，则 G 是可除的。

事实上，存在一个紧连通群 K 及一个整数 \(n \geqslant 0\)，使得 G 同构于 \(R^{n} \times K\)（II，第 248 页，命题 3，其中必有 L = K）。推论 1 表明 K 是可除的，因而 G 是可除的。

